\documentclass[12pt]{article}
\usepackage{amsfonts,amsthm,amsmath,amssymb,mathtools}
\usepackage{mathrsfs}
\usepackage{enumitem}
\usepackage{tikz-cd}
\usepackage{geometry}
\usepackage{mlmodern}
\usepackage{graphicx}

\geometry{letterpaper, portrait, margin=1in}

\setcounter{section}{0}

\renewcommand{\thesection}{\S\arabic{section}}
\renewcommand{\thesubsection}{\arabic{section}}
\DeclareMathOperator{\im}{im}
\DeclareMathOperator{\id}{id}
\DeclareMathOperator{\ad}{ad}
\newcommand{\gl}{\mathfrak{gl}}
\renewcommand{\sl}{\mathfrak{sl}}
\renewcommand{\th}{^\mathrm{th}}

\newtheorem{thm}{Theorem}
\newenvironment{theorem}[1]{
	\renewcommand{\thethm}{#1}
	\begin{thm}
		}{\end{thm}}

\newtheorem{lma}{Lemma}
\newenvironment{lemma}[1]{
	\renewcommand{\thelma}{#1}
	\begin{lma}
		}{\end{lma}}

\theoremstyle{definition}
\newtheorem{ex}{}
\newenvironment{exercise}[1]{
	\renewcommand{\theex}{#1}
	\begin{ex}
		}{\end{ex}}

\title{Lie Theory: Chapter 1 Homework}
\author{Connor Mooneyhan}
\date{}

\begin{document}

\maketitle

\begin{ex}
	Let $\mathbb{R}^3_\vee$ be the Lie algebra defined in Section 1.2.
	Prove that if $x \cdot y$ denotes the dot product of vectors $x, y \in \mathbb{R}^3$, then \begin{equation*}
		x \wedge (y \wedge z) = (x \cdot z) y - (x \cdot y) z \hspace{1em}\forall x, y \in \mathbb{R}^3.
	\end{equation*}
\end{ex}
\begin{proof}
	Observe that
	\begin{align*}
		x \wedge (y \wedge z) =  ( & x_1, x_2, x_3) \wedge (y_2z_3 - y_3z_2, y_3z_1 - y_1z_3, y_1z_2 - y_2z_1) \\
		=                        ( & x_2y_1z_2 - x_2y_2z_1 - x_3y_3z_1 + x_3y_1z_3,                            \\
		                           & x_3y_2z_3 - x_3y_3z_2 - x_1y_1z_2 + x_1y_2z_1,                            \\
		                           & x_1y_3z_1 - x_1y_1z_3 - x_2y_2z_3 + x_2y_3z_2)                            \\
		=                        ( & x_2y_1z_2 - x_2y_2z_1 - x_3y_3z_1 + x_3y_1z_3,                            \\
		                           & x_3y_2z_3 - x_3y_3z_2 - x_1y_1z_2 + x_1y_2z_1,                            \\
		                           & x_1y_3z_1 - x_1y_1z_3 - x_2y_2z_3 + x_2y_3z_2)                            \\
	\end{align*}
	and \begin{align*}
		= (                        & x \cdot z)y - (x \cdot y)z                                                                               \\
		= (                        & x_1z_1 + x_2z_2 + x_3z_3)(y_1, y_2, y_3) - (x_1y_1 + x_2y_2 + x_3y_3)(z_1, z_2, z_3)                     \\
		= (                        & x_1z_1y_1 + x_2z_2y_1 + x_3z_3y_1, x_1z_1y_2 + x_2z_2y_2 + x_3z_3y_2, x_1z_1y_3 + x_2z_2y_3 + x_3z_3y_3) \\
		- (                        & x_1y_1z_1 + x_2y_2z_1 + x_3y_3z_1, x_1y_1z_2 + x_2y_2z_2 + x_3y_3z_2, x_1y_1z_3 + x_2y_2z_3 + x_3y_3z_3) \\
		=                        ( & x_2y_1z_2 - x_2y_2z_1 - x_3y_3z_1 + x_3y_1z_3,                                                           \\
		                           & x_3y_2z_3 - x_3y_3z_2 - x_1y_1z_2 + x_1y_2z_1,                                                           \\
		                           & x_1y_3z_1 - x_1y_1z_3 - x_2y_2z_3 + x_2y_3z_2),
	\end{align*}
	where the the last line comes after commuting elements in the terms with same index across all three variables (i.e. $x_1z_1y_1$ becomes $x_1y_1z_1$) since $\mathbb{R}$ is commutative and then cancelling them out.

	We can see that both expressions reduce to the same final expression, so they are equal.
\end{proof}

\begin{ex}
	Check that \begin{equation*}
		[e_{ij}, e_{kl}] = \delta_{jk}e_{il}-\delta_{il}e_{kj},
	\end{equation*}
	where $\delta$ is the Kronecker delta, defined by $\delta_{ij} = 1$ if $i = j$ and $\delta_{ij} = 0$ otherwise.
\end{ex}
\begin{proof}
	Recall that $[e_{ij}, e_{kl}] = e_{ij}e_{kl} - e_{kl}e_{ij}$.
	The first term will have zeroes in all but the $i\th$ row and zeroes in all but the $l\th$ column since the only nonzero row in $e_{ij}$ is the $i\th$ one and the only nonzero column in $e_{kl}$ is the $l\th$ one.
	Thus the only possibly-nonzero entry is the one in the $i\th$ row and $l\th$ column.
	This will be 1 if and only if $j = k$ (and 0 otherwise) since we will then have the 1 from $e_{ij}$ multiplied by the 1 in $e_{kl}$ and zeroes everywhere else when doing the matrix multiplication. Thus we obtain the first term $\delta_{jk}e_{il}$ of our desired result.

	The only other term to deal with is the second one, for which we use the same argumentation on $-e_{kl}e_{ij}$ to conclude that we subtract a 1 from the entry at the $j\th$ column of the $k\th$ row if and only if $i = l$, and a zero otherwise.
	This gives us our second term of $-\delta_{il}e_{kj}$, and we are done.
\end{proof}

\begin{ex}
	Prove that the set $e_{ij}$ for $i \neq j$ together with the set $e_{ii} - e_{i+1,i+1}$ for $1 \leq i < n$ forms a basis for $\sl(n, F)$.
\end{ex}

\begin{ex}
	Prove that an ideal of a Lie algebra is always a subalgebra.
\end{ex}
\begin{proof}
	Let $I$ be an ideal of a Lie algebra $L$.
	Then $I$ is a linear subspace of $L$ by virtue of being an ideal, so we just need to show that $[x, y] \in I$ whenever $x, y \in I$.
	This is quite trivial, since $I$ being an ideal guarantees that for any $x \in L$ and $y \in I$, $[x, y] \in I$, and surely $x \in L$ if $x \in I \subset L$, so there is nothing more to show.
\end{proof}

\begin{ex}
	Consider the adjoint homomorphism map defined (on page 4) for all $x, y \in L$ by \begin{equation*}
		(\ad x)(y) \coloneq [x, y].
	\end{equation*}
	\begin{enumerate}[label=(\alph*)]
		\item Prove that $\ad$ is linear.
		\item Prove that $\ad$ is a homomorphism.
		\item Prove that the kernel of $\ad$ is the center of $L$.
	\end{enumerate}
\end{ex}
\begin{proof}\
	\begin{enumerate}[label=(\alph*)]
		\item For $a, b \in F$ and $x, y, z \in L$, \begin{align*}
			      (\ad (ax + bz)(y)) & = [ax + bz, y]                                                          \\
			                         & = a[x, y] + b[z, y] \hspace{1em} \text{(by bilinearity of the bracket)} \\
			                         & = a(\ad x)(y) + b(\ad z)(y).
		      \end{align*}
		\item We have already shown that $\ad$ is linear, so we just need to show that $\ad[x, y] = [\ad x, \ad y]$.
		      Indeed, for any $x, y, z \in L$ \begin{align*}
			      (\ad[x, y])(z) & = [[x, y], z]                                             \\
			                     & = - [z, [x, y]]                                           \\
			                     & = [x, [y, z]] + [y, [z, x]]                               \\
			                     & = [x, [y, z]] -  [y, [x, z]]                              \\
			                     & = (\ad x) ([y, z]) - (\ad y) ([x, z])                     \\
			                     & = (\ad x) ((\ad y)(z)) - (\ad y) ((\ad x)(z))             \\
			                     & = ((\ad x) \circ (\ad y))(z) - ((\ad y) \circ (\ad x))(z) \\
			                     & = [\ad x, \ad y](z).
		      \end{align*}
		\item The center of $L$ is defined as the set of elements $x$ of $L$ that satisfy $[x, y] = 0$ for any $y \in L$.
		      If $x \in L$ is in the kernel of $\ad$, then $\ad x$ is the zero map, so by definition $0 = (\ad x)(y) = [x, y]$ for all $y \in L$.
		      Thus elements of the kernel of $\ad$ are in the center, and just as easily we see that if $x$ is in the center of $L$, $[x, y] = 0$ for all $y \in L$, so $\ad x$ is the zero map.
		      By mutual inclusion then, the kernel of $\ad$ is the center of $L$.
	\end{enumerate}
\end{proof}

\begin{ex}
	Prove that $\mathrm{Der}(A)$ is a vector subspace of $\gl(A)$.
\end{ex}
\begin{proof}
	By definition, derivations are linear, so $\mathrm{Der}(A) \subset \gl(A)$.
	We need to show, then, that given an algebra $A$ over a field $F$ along with derivations $D, E \in \mathrm{Der}(A)$ and $c \in F$, $D + E \in \mathrm{Der}(A)$ and $cD \in \mathrm{Der}(A)$.
	Indeed, for any $a, b \in A$, we have \begin{align*}
		(D + E)(ab) & = D(ab) + E(ab)                   \\
		            & = aD(b) + D(a)b + aE(b) + E(a)b   \\
		            & = aD(b) + aE(b) + D(a)b + E(a)b   \\
		            & = a(D(b) + E(b)) + (D(a) + E(a))b \\
		            & = a(D + E)(b) + (D + E)(a)b
	\end{align*}
	and \begin{align*}
		(cD)(ab) & = c(D(ab))                                                                           \\
		         & = c(aD(b) + D(a)b)                                                                   \\
		         & = c(aD(b)) + c(D(a)b)                                                                \\
		         & = a(cD(b)) + (cD(a))b \hspace{1em} \text{(by bilinearity of algebra multiplication)} \\
		         & = a(cD)(b) + (cD)(a)b,
	\end{align*}
	so both satisfy the Leibniz rule and are thus derivations as desired.
\end{proof}

\begin{ex}
	Let $L_1$ and $L_2$ be two abelian Lie algebras.
	Show that $L_1$ and $L_2$ are isomorphic if and only if they have the same dimension.
\end{ex}
\begin{proof}
	Suppose $L_1$ and $L_2$ both have dimension $n$, and let $\lbrace x_1, \dots, x_n \rbrace$ and $\lbrace y_1, \dots, y_n \rbrace$ be bases of $L_1$ and $L_2$ respectively.
	Then we can define a function $\varphi : L_1 \to L_2$ by letting $\varphi(x_i) = y_i$ for each $i$, and \begin{align*}
		\varphi(a_1x_1 + \cdots + a_nx_n) & \coloneq a_1\varphi(x_1) + \cdots + a_n\varphi(x_n) \\
		                                  & = a_1y_1 + \cdots + a_ny_n.
	\end{align*}
	This is well-defined since every element of $L_1$ can be expressed uniquely as a linear combination of $\lbrace x_1, \dots, x_n \rbrace$.
	It is also linear by construction, it is injective since linear combinations of $\lbrace y_1, \dots, y_n \rbrace$ are unique, and it is surjective because \textit{every} element of $L_2$ can be expressed as such a linear combination.
	That leaves us to confirm that $\varphi([x, y]) = [\varphi(x), \varphi(y)]$.
	Indeed, given $x, y \in L_1$, \begin{align*}
		    & [\varphi(a_1x_1 + \cdots + a_nx_n), \varphi(b_1x_1 + \cdots + b_nx_n)]                                                                     \\
		=\  & [a_1y_1 + \cdots + a_ny_n, b_1y_1 + \cdots + b_ny_n]                                                                                       \\
		=\  & a_1[y_1, b_1y_1 + \cdots + b_ny_n] + \cdots + a_n[y_n, b_1y_1 + \cdots + b_ny_n]                                                           \\
		=\  & (a_1b_1[y_1, y_1] + a_1b_2[y_1, y_2] + \cdots + a_1b_n[y_1, y_n]) + \cdots + (a_nb_1[y_n, y_1] + a_nb_2[y_n, y_2] + \cdots + a_nb_n[y_n, y_n]) \\
		=\  & \cdots                                                                                                                                     \\
		=\  & \varphi([a_1x_1 + \cdots + a_nx_n, b_1x_1 + \cdots + b_nx_n])                                                                              \\
		=\  & \varphi([a_1x_1 + \cdots + a_nx_n, b_1x_1 + \cdots + b_nx_n])                                                                              \\
	\end{align*}
\end{proof}

\begin{ex}
	Find the structure constants of $\sl(2, F)$ with respect to the basis given by the matrices \begin{equation*}
		e = \begin{pmatrix} 0 & 1 \\ 0 & 0 \end{pmatrix}, \hspace{0.25em} f = \begin{pmatrix}0 & 0 \\ 1 & 0\end{pmatrix}, \hspace{0.25em} h = \begin{pmatrix}1 & 0 \\ 0 & -1\end{pmatrix}.
	\end{equation*}
\end{ex}

\begin{ex}
	Prove that $\sl(2, \mathbb{C})$ has no non-trivial ideals.
\end{ex}

\begin{ex}
	Let $S$ be an $n \times n$ matrix with entries in a field $F$.
	Define \begin{equation*}
		\gl_S(n, F) = \left\lbrace x \in \gl(n, F) : x^T S = -S x \right\rbrace.
	\end{equation*}
	\begin{enumerate}[label=(\alph*)]
		\item Show that $\gl_S(n, F)$ is a Lie subalgebra of $\gl(n, F)$.
		\item Find $\gl_S(2, \mathbb{R})$ if $S = \begin{pmatrix}0 & 1\\0&0\end{pmatrix}$.
	\end{enumerate}
\end{ex}

\begin{ex}
	Let $A$ be an algebra and let $\delta : A \to A $ be a derivation.
	Prove that $\delta$ satisfies the Leibniz Rule \begin{equation*}
		\delta^n(xy) = \sum_{r=0}^n \begin{pmatrix}n\\r\end{pmatrix} \delta^r(x) \delta^{n-r}(y) \hspace{1em} \forall x, y \in A.
	\end{equation*}
\end{ex}

\end{document}
